\label{ch:payment-scope}
\section{Current Assessment Decision}
\begin{decisionbox}[Real payment processing is intentionally not implemented]
The submitted backend focuses on finite ticket inventory, retry safety,
voucher consistency, booking lifecycle, and flash-sale behavior. A newly
created booking is a time-limited reservation. Confirmation is simulated
through the owned customer command or an authorized operation workflow; it does
not prove card, bank, wallet, or provider settlement.
\end{decisionbox}

Adding a provider SDK would introduce network failure, webhook authentication,
settlement, refunds, reconciliation, and sensitive-data concerns without
materially strengthening evidence for the assessment's core correctness risks.
The omission is therefore a deliberate scope decision, not an accidental gap.

\section{Implemented Payment-Adjacent Behavior}
The current code supports the parts of booking lifecycle that a later payment
orchestration module would depend on:
\begin{itemize}
  \item create a ten-minute \code{RESERVED} booking and atomically hold tickets and optional voucher capacity;
  \item persist booking state, immutable totals, expiry time, and audit history;
  \item conditionally transition a live \code{RESERVED} booking to \code{CONFIRMED};
  \item expire or cancel a reservation and release finite resources exactly once;
  \item reject a terminal-state race with a stable conflict;
  \item apply idempotency to booking creation.
\end{itemize}

\section{Explicitly Not Implemented}
There is no external payment-provider client, payment session, authorization or
capture, card/wallet/bank handling, provider webhook verification, payment
status table, settlement, refund, chargeback, reconciliation, PCI-DSS card-data
handling, or payment-specific idempotency workflow. The backend cannot accept
or verify real customer payments.

\section{Why Production Uses an External Gateway}
A production system normally delegates sensitive payment execution to a
specialized external gateway. The backend still owns payment orchestration,
the booking/payment association, persisted payment state, verified and
idempotent callback processing, booking transitions, expiry rules, and
reconciliation jobs. The provider owns the customer payment interaction and
regulated processing surface.

Figure~\ref{fig:future-payment-flow} is an evolution model only; none of its
dashed provider interactions exists in the submitted code.

\begin{figure}[H]
\centering
\begin{tikzpicture}[
  node distance=0.95cm,
  current/.style={draw=Primary,rounded corners,fill=SoftBlue,minimum width=10cm,minimum height=0.70cm,align=center},
  future/.style={draw=gray,rounded corners,fill=SoftGray,dashed,minimum width=10cm,minimum height=0.70cm,align=center},
  arrow/.style={-{Latex[length=2mm]},thick,draw=Primary},
  futurearrow/.style={-{Latex[length=2mm]},thick,dashed,draw=gray}
]
\node[current] (reserve) {Transaction 1: reserve resources and commit RESERVED booking};
\node[future,below=of reserve] (session) {Future: request hosted payment session after commit};
\node[future,below=of session] (customerpay) {Future: customer completes payment at external gateway};
\node[future,below=of customerpay] (webhook) {Future: verify and idempotently persist provider webhook};
\node[current,below=of webhook] (confirm) {Transaction 2: conditionally transition live RESERVED to CONFIRMED};
\draw[futurearrow] (reserve) -- (session);
\draw[futurearrow] (session) -- (customerpay);
\draw[futurearrow] (customerpay) -- (webhook);
\draw[futurearrow] (webhook) -- (confirm);
\end{tikzpicture}
\caption{Future external payment flow (not implemented)}
\label{fig:future-payment-flow}
\end{figure}

If payment completes after expiry wins, the backend must not resurrect the
terminal booking. A production policy would record the late provider result
and start a refund/compensation or manual reconciliation path.

\section{Transaction Design Rule}
The ticket reservation transaction must never remain open while waiting on a
provider network call. External systems cannot participate in the same local
ACID transaction; holding row locks across network latency amplifies hot-row
contention and makes provider timeouts database failures.

\begin{table}[H]
\centering
\caption{Payment transaction boundary}
\begin{tabularx}{\textwidth}{p{0.20\textwidth}YY}
\toprule
\textbf{Phase} & \textbf{Preferred} & \textbf{Avoid} \\
\midrule
Reservation & Reserve inventory, create \code{RESERVED}, commit & Keep transaction open while calling provider \\
Provider interaction & Occur after commit with its own timeout/idempotency & Treat network response as part of local ACID \\
Result handling & Verify callback; begin a new transaction; conditionally confirm & Blindly overwrite an expired/cancelled booking \\
Failure recovery & TTL release plus reconciliation/compensation & Indefinite resource hold or untracked dual write \\
\bottomrule
\end{tabularx}
\end{table}

\section{Limitation and Evolution Statement}
The implementation does not address webhook signature/replay security,
provider outage and retry policy, settlement mismatch, refund/chargeback
workflow, or financial reconciliation. A production evolution should add a
separate payment orchestration module and payment-attempt model, provider
adapter, verified idempotent webhook inbox, explicit payment state machine,
transactional outbox for side effects, reconciliation jobs, and operational
alerts---without changing PostgreSQL's authority over booking state.
